% TCC — Desenvolvimento
\section{Desenvolvimento da Ferramenta}\label{sec:desenvolvimento}

\subsection{Arquitetura e Fluxo Operacional}

O OpenAPI Contract Fuzzer foi implementado em Python como um pipeline modular.
O contrato é convertido em um modelo interno; esse modelo orienta a geração de
valores-base e mutações; um plano materializa as requisições; o executor envia
as chamadas à API; e o analisador confronta as respostas com o contrato antes
de encaminhar os resultados reportáveis ao sistema fuzzy e ao gerador de
relatórios.

% TODO(TCC): inserir diagrama editável:
% OpenAPI -> Parser -> Modelo -> Geração/Mutação -> Plano/Runner -> API
% -> Analisador Contratual -> FIS -> Relatório.

\subsection{Interpretação do Contrato e Geração de Entradas}

O parser atual aceita documentos YAML ou JSON e representa operações,
parâmetros, corpos JSON, respostas e mecanismos de segurança dentro de um
subconjunto de OpenAPI. A partir do modelo, o gerador produz valores-base e o
motor aplica operações de substituição ou remoção sobre tipos primitivos,
objetos e arrays.

O suporte não é integral: referências e composições de schema, diferentes
\textit{media types} e parte das regras de serialização ainda não são
interpretados. Além disso, o catálogo existente é predominantemente orientado
por tipos. O texto e a avaliação deverão distinguir restrições efetivamente
usadas de capacidades ainda não implementadas.

\subsection{Planejamento e Execução}

O plano de execução associa operação, caminho resolvido, corpo e cabeçalhos.
Autenticação pode ser fornecida por configuração ou obtida por um procedimento
automático baseado no contrato. A execução utiliza um cliente HTTP assíncrono e
um semáforo para limitar a concorrência, preservando a ordem lógica entre plano
e resultados.

\subsection{Detecção de Divergências}

O analisador identifica respostas sem status, erros de servidor, status não
documentados, incompatibilidades de schema e aceitação de entradas mutadas.
Uma rejeição 4xx documentada e compatível com o schema é comportamento esperado
e não deve ser contabilizada como anomalia candidata.

Um achado reportado não equivale automaticamente a defeito ou vulnerabilidade.
A confirmação exige reprodução e inspeção das condições do contrato, da entrada
e do ambiente.

\subsection{Priorização por Lógica Fuzzy}

Após a classificação contratual, o sistema Mamdani combina desvio de status,
tempo de resposta e tamanho da resposta para produzir um escore contínuo,
posteriormente dividido nas faixas LOW, MEDIUM e HIGH. As funções de
pertinência, regras e limiares atuais são escolhas heurísticas do autor; sua
coerência constitui objeto da avaliação, não uma propriedade presumida.

\subsection{Artefatos e Rastreabilidade}

Cada campanha produz plano, resultados brutos, achados estruturados e relatório
em Markdown. Antes da coleta definitiva, os artefatos deverão incluir
identificador de caso, mutação aplicada, versão do contrato, configuração e
seed, permitindo deduplicação, reexecução e rastreabilidade dos números
apresentados no artigo.
